\documentclass[11pt]{article}
\usepackage[a4paper,margin=1in]{geometry}
\usepackage{amsmath,amssymb}
\usepackage{booktabs}
\usepackage{array}
\usepackage{listings}
\usepackage{hyperref}
\usepackage{xcolor}

\lstset{basicstyle=\ttfamily\small,breaklines=true,frame=single,framerule=0.2pt,
        columns=fullflexible,keepspaces=true}

\title{The Scan Algebra $\mathfrak{S}$: A Unified Derivation from Grammars to Parallel Parsers}
\author{pars}
\date{\today}

\begin{document}
\maketitle

\begin{abstract}
We formalize the parallelism behind cuJSON, simdjson and simdcsv as a single
the scan algebra $\mathfrak{S}$: a small algebra of primitives (classification, monoid scan, masked
select, compaction, well-nested reduction, associative reduction) together with
a rewrite system that turns a naive branch program into branchless masked data
flow.  We show that a grammar lowers to $\mathfrak{S}$ exactly when its rule bodies are
regular and its recursion is well-nested or associative, that cuJSON/simdjson/
simdcsv are three instances of $\mathfrak{S}$, and that the resulting pipelines are
generated without any format-specific code.
\end{abstract}

\section{Naive branch programs}
A \emph{naive branch program} (NBP) operates positionally on a byte stream using
only: \texttt{if/else} whose conditions are the current byte class and a bounded
register state; loops over positions; bounded registers; a stack; and
\texttt{emit}.  The whole of cuJSON's tokenizer is the NBP
\begin{lstlisting}
in_string = escaped = false;
for (i = 0; i < n; ++i) { c = data[i];
  if (in_string) { if (escaped) escaped=false;
                   else if (c=='\\') escaped=true;
                   else if (c=='"') in_string=false; }
  else { if (c=='"') in_string=true;
         else if (is_structural(c)) { emit_external(i);
              if (is_bracket(c)) emit_bracket(i); } } }
\end{lstlisting}
and its parser is the stack loop
\begin{lstlisting}
for (tok in bracket_indices)
  if (is_open(data[tok])) push(tok);
  else { j = pop(); pair[j] = tok; }
\end{lstlisting}

\section{The rewrite system}
Let $\mathrm{class}:\Sigma\to C$ be derived from the terminal sets.

\begin{description}
\item[R1 (classification)] $c=\ell\mapsto M_\ell$; $c\in S\mapsto\bigvee_{s\in S}M_s$.
\item[R2 (state $\to$ mask)] a $k$-bit register becomes $k$ bit masks.
\item[R3 (branch $\to$ select)]
$\texttt{if}(c)\,x{=}f(x)\,\texttt{else}\,x{=}g(x)\;\rightsquigarrow\;
x=(c\wedge f(x))\vee(\neg c\wedge g(x))$.
\item[R4 (scan)] the recurrence $s_{i+1}=T_{\mathrm{class}_i}(s_i)$ makes each byte
induce a transform $\tau_i:S\to S$, so
\[
s_i=(\tau_{i-1}\circ\cdots\circ\tau_0)(s_0),
\]
a prefix product in the transformation monoid $(S\to S,\circ)$, computed by a
parallel scan.  For $S=\{0,1\}$ the transform is affine,
\[
\tau(x)=(a\wedge x)\oplus b,\qquad
(a_1,b_1)\circ(a_2,b_2)=(a_1\wedge a_2,\,(b_1\wedge a_2)\oplus b_2),
\]
whose prefix is simdjson's \texttt{prefix\_xor} and cuJSON's escape scan.
\item[R5 (compaction)] $\texttt{for}\,i\,\texttt{if}\,M[i]\,\texttt{out.push}(i)
\rightsquigarrow \mathrm{popcount}+\mathrm{exclusive\_scan}+\mathrm{scatter}$.
\item[R6 (well-nested)] for a well-nested bracket language the stack content is
determined by the depth $\delta_i$, the prefix sum of $(+1,-1)$ (a scan over the
group $\mathbb Z$); stable-sort by depth and pair adjacent tokens, then validate
types.  This is cuJSON's \texttt{depth}$\to$\texttt{sort}$\to$\texttt{validate}.
\item[R7 (associative)] a single self-recursive rule with an associative operator
lowers to a monoid reduction.
\end{description}

\begin{definition}[$\mathfrak{S}$-reducibility]
A grammar $G$ is \emph{$\mathfrak{S}$-reducible} iff
(i) every non-recursive rule body is a regular expression over the terminal
alphabet, and
(ii) every recursive strongly connected component of the rule graph is either
\emph{visibly pushdown} (each recursive call is guarded by a terminal
call/return pair) or a single self-recursive rule whose operator is associative.
Condition (i) yields R1--R5; condition (ii) yields R6 or R7.
\end{definition}

\section{Three SOTA parsers are instances of $\mathfrak{S}$}
\begin{center}
\begin{tabular}{@{}ll@{}}
\toprule
Parser & $\mathfrak{S}$ decomposition\\
\midrule
simdjson stage1 & R1 + R4(\texttt{prefix\_xor},\texttt{escaped\_scan}) + R3 + R5\\
cuJSON tokenize & R1 + R4 ($\texttt{\_\_vcmpeq4}$) + R3 + R5\\
cuJSON parser & R6 (depth $=$ R4 over $\mathbb Z$; pair $=$ sort by depth)\\
simdcsv & R1 + R4 (quote parity) + R3 + R5\\
\bottomrule
\end{tabular}
\end{center}

\section{Compilation (shape-free)}
\texttt{make\_scan\_plan}$\langle G\rangle$ uses no syntactic shape patterns.
\begin{enumerate}
\item \textbf{Lexical rules $\to$ DFA.}  Reify the rule body to a regular
expression and build a DFA by Brzozowski derivatives; states are canonical
derivatives, $\mathsf{alt}$ is normalized by ACI sorting/dedup, and the alphabet
is the quotient of bytes by membership in every character predicate (so a dense
complement predicate still yields a tiny alphabet).
\item \textbf{Masks from the automaton.}  A byte is masked iff the DFA state
before it is non-start and can reach accept through a path consuming at least
one structural byte.  One definition covers prefix-escape, no-escape, doubled
delimiters and line comments.  The \emph{scheduled} delimiter/escape form is read
off the DFA: a unique start byte $D$; $\mathsf{der}(S_1,D)$ with
$\mathsf{der}(S_2,D)=S_1$ gives a doubled delimiter, an absorbing return through
some $E$ gives a prefix escape, otherwise no escape.
\item \textbf{Brackets from FIRST/LAST.}  For a recursive SCC whose FIRST and
LAST terminal sets are single distinct characters, emit an R6 bracket pair ---
the visibly-pushdown guard, not a syntactic position.
\item \textbf{Punctuation} is the standalone terminals of non-lexical rules.
\item \textbf{Support.}  Every recursive rule must be R6-guarded; else diagnose.
\end{enumerate}
The lexical/syntactic split is the standard BNF annotation \texttt{lexical\_rules}
(inferred as non-recursive rules when absent).

\paragraph{Execution semantics (P1).}
The plan carries, per lexical masking rule, a monoid descriptor (still an NTTP):
\texttt{affine} (delimiter set $\times$ optional prefix escape, executed as
$x'{=}(a\wedge x)\oplus b$), \texttt{doubled}, \texttt{comment}
(start set $+$ newline reset), or \texttt{table} (a $\le 8\times 8$ transition
table).  Execution is a per-byte state transform whose prefix product is the
running state (R4).  The mask is \emph{not} a per-state predicate but the
per-transition bit ``this byte is consumed by an in-progress token'', which is
exact for doubled delimiters (a per-state inside bit is off by one at the
closing quote).  Factorization is read off the DFA and validated by differential
random testing.  The decisive coverage witness is a block comment
$/*\cdots*/$ (multi-character, two distinct delimiters), which is not one of the
legacy templates and runs on the \texttt{table} path; the GPU consuming the same
table (P1c) and a formal specialization-equivalence proof (P1d) remain.

\paragraph{Soundness gates.}
The unified mask predicate is executable: a lexical rule becomes a mode/comment
only if some non-start DFA state can reach accept through a structural byte.
General \texttt{notp} is rejected (only the one-byte complement form
$\texttt{seq}{<}\texttt{notp}{<}\texttt{cls}{<}S{>}{>},\texttt{any}{>}$ is
regular).  A recursive rule is accepted as an R6 anchor only if it satisfies the
conservative structural criterion ``the top-level sequence has exactly the two
literal endpoints, distinct, with the recursion strictly inside'', which rejects
ill-founded patterns such as $a \to (\,a\,)\,;$; well-nestedness is additionally
validated at run time.  \texttt{require\_supported} makes any unsupported rule a
compile error unless it is listed in \texttt{fallback\_rules}.

\section{Implementation and results}
The framework contains no format-specific code; grammars are declared on the
bench side only.  Throughput on 67.1\,MB payloads (RTX 4060 Laptop, CUDA 13.4):
\begin{center}
\begin{tabular}{@{}llr@{}}
\toprule
Format & Implementation & GB/s\\
\midrule
JSON & pars $\mathfrak{S}$ + CUDA & 0.98--1.23\\
JSON & cuJSON (hand-written CUDA) & $\sim$0.96--1.1\\
JSON & simdjson (CPU ondemand) & 1.49\\
CSV  & pars $\mathfrak{S}$ + CUDA & 2.15\\
CSV  & simdcsv (CPU AVX2) & 9.37\\
\bottomrule
\end{tabular}
\end{center}
JSON matches cuJSON's order of magnitude.  CSV trails the hand-written AVX2
parser, mainly because the host-to-device copy is not amortized at 67\,MB and
CSV has low structural density.  Our implementation also corrects a real bug in
the adjacency-based doubled-quote detection (it under-counts); parity toggling
matches the naive reference ($8{,}567{,}094$ separators).

\section{Cost model and the CSV anomaly}
The primitives above characterize \emph{which computations} are data-parallel;
they do not by themselves yield end-to-end throughput.  Make the schedule part of
the model.  For an $n$-byte input producing $k$ structural tokens,
\[
T \;=\; T_{\mathrm{move\text{-}in}}(n) \;+\; T_{\mathrm{compute}}(n,k)
\;+\; T_{\mathrm{move\text{-}out}}(k) \;+\; T_{\mathrm{intermediate}} .
\]
\begin{theorem}[Movement bound]
R1--R7 give $T_{\mathrm{compute}}(n,k)=O(n\,\mathrm{polylog}\,n / \mathrm{BW}_{\mathrm{gpu}})$
with $O(\log n)$ depth, independent of $T_{\mathrm{move}}$.  Hence end-to-end
throughput $n/T$ is governed by $\max(T_{\mathrm{compute}},T_{\mathrm{move}})$; if
primitives materialize $\Theta(n)$ intermediates and results are copied to the
host, $T_{\mathrm{move}}$ dominates and the parallelism of $\mathfrak{S}$ is unobservable.
\end{theorem}

\paragraph{Measured attribution (67.1\,MB, RTX 4060 Laptop).}
\begin{center}
\begin{tabular}{@{}lrr@{}}
\toprule
Configuration & JSON (GB/s) & CSV (GB/s)\\
\midrule
kernels only (device-resident) & 62.5 & 71.6\\
$+$ H2D input copy & 3.81 & 6.61\\
$+$ D2H result copy-back (structural only) & 1.78 & 3.91\\
\bottomrule
\end{tabular}
\end{center}
$T_{\mathrm{compute}}$ is $60$--$70$\,GB/s; the end-to-end number is $1$--$3$\,GB/s.
\texttt{nsys} attributes $63.7\%$ of memory-op time to D2H (full structural and
pair arrays) and $35.1\%$ to H2D (pageable $7.9$\,GB/s; pinned $12.4$\,GB/s).
CSV exposes this sharply: its structural density is low ($6/47\approx13\%$) and it
has no nesting, so $T_{\mathrm{compute}}$ is nearly free and the fixed
$T_{\mathrm{move}}$ is the entire cost.  This is \emph{not} an algorithmic loss;
it is a placement/materialization loss.  Consequently $\mathfrak{S}$ is extended with a
schedule pass (placement, fusion, streaming, and an explicit observable
boundary), and throughput claims must distinguish device-resident
$T_{\mathrm{compute}}$ from end-to-end $n/T$.

\section{Related work and positioning}
The algebraic components of $\mathfrak{S}$ are classical: transformation monoids and
parallel prefix are Blelloch's scan; automata-as-monoids is textbook; the
well-nested reduction is cuJSON (ASPLOS'26); the bracket/stack monoid is Levien's
\emph{Stack Monoid} (2020); well-nested recursion is precisely visibly pushdown
(Alur--Madhusudan, STOC'04).  Our claim is therefore \emph{not} the algebra but
its \emph{systematization and compilation}: a mechanical derivation
\textsf{grammar}$\to$\textsf{primitives}$\to$\textsf{schedule}, with the
compute/movement distinction made explicit.  This compiler view is absent from
the prior systems (simdjson, simdcsv, cuJSON, GpJSON, Mison, Sparser, Pison).

\section{Boundary}
$\mathfrak{S}$ covers grammars with regular rule bodies and well-nested or associative
recursion.  Unbounded lookahead and non-well-nested recursion are outside $\mathfrak{S}$ and
must be declared as fallbacks; stage-2 semantic parsing is sequential/on-demand
and is likewise outside $\mathfrak{S}$.

\begin{thebibliography}{9}
\bibitem{blelloch} G.~E. Blelloch. \emph{Prefix sums and their applications}.
  Technical Report CMU-CS-90-190, 1990.
\bibitem{simdjson} G.~Langdale and D.~Lemire. Parsing gigabytes of JSON per
  second. \emph{VLDB Journal}, 2019.
\bibitem{cujson} A.~Vedadi Gargary, S.~Safari Loaliyan, Z.~Zhao. cuJSON: A Highly
  Parallel JSON Parser for GPUs. \emph{ASPLOS}, 2026.
\bibitem{simdcsv} G.~Langdale. simdcsv, 2019.
\bibitem{gpjson} GpJSON: GPU-accelerated JSON parsing. \emph{VLDB}, 2025.
\bibitem{mison} L.~Li et al. Mison: A Fast JSON Parser for Data Analytics.
  \emph{VLDB}, 2017.
\bibitem{sparser} S.~Palkar et al. Filter Before You Parse: Faster Analytics on
  Raw Data. \emph{VLDB}, 2018.
\bibitem{pison} L.~Li et al. Pison: A Parallel JSON Parser for Many-core.
\bibitem{alur} R.~Alur and P.~Madhusudan. Visibly Pushdown Languages.
  \emph{STOC}, 2004.
\bibitem{stackmonoid} R.~Levien. The Stack Monoid, 2020.
\end{thebibliography}

\end{document}
